\chapter*{Prefacio}

Este es un curso avanzado sobre dinámica estructural, que incluye aplicaciones en el modelamiento de sistemas dinámicos sometidos a cargas estocásticas, sistemas de monitoreo y control estructural, problemas inversos o de identificación de sistemas, entre otros. Los conceptos y las técnicas cubiertas en este curso son fundamentales para realizar investigación aplicada y desarrollo tecnológico en el campo de la ingeniería estructural, con énfasis en el diseño de sistemas estructurales inteligentes (\emph{smart structures}).

El apunte se divide en tres grandes partes. La primera parte se enfoca principalmente en establecer las bases de la teoría de sistemas dinámicos lineales contínuos y discretos en el tiempo, sujetos a cargas determinísticas o estocásticas. La segunda parte, se concentra en el desarrollo de aplicaciones avanzadas en dinámica estructural, principalmente lo relacionado con monitoreo y control estructural. Finalmente, la tercera parte esta dedicado a tópicos avanzados que avanzan muy rápidamente junto con el estado del arte en dinámica estructural. Temas como simulación híbrida y aprendizaje automático son sólo una selección de tópicos cubiertos en esta última parte.

Agradecimientos a María Elena Quiroz (M.S., Ing. USM), quien fue la segunda ayudante de docencia del curso IPO426 en el primer semestre de 2022. María Elena fue quien tomó mis apuntes tomados a mano en las primeras instancias del curso, y las transcribió en este documento. Además, fue ella la que preparó la totalidad de figuras e ilustraciones del primer borrador de este apunte de clases. También agradezco a Alexis Contreras (Ing. y estudiante magíster USM), el tercer ayudante de docencia del curso IPO426/MIC411 realizado el primer semestre de 2024. En particular, Alexis ayudó a revisar todo el contenido de estos apuntes, y también colaboró en la elaboración de nuevos capítulos, como por ejemplo el de sensores. Por último, agradecimientos especiales a Cristobal Gálmez (M.S., Ing. USM), el primer ayudante de docencia del curso IPO426 realizado el primer semestre de 2020.

\bigskip\noindent
\begin{flushright}
  Gastón Fermandois\\
  5 de agosto de 2025\\
  Valparaíso, Chile
\end{flushright}